% Source: Weinberg GR, physical PDF pp. 322--327
%         (printed pp. 299--304).
% Coverage: Section 11.1, Differential Equations for Stellar Structure;
%           equations (11.1.1)--(11.1.31).
% Figures/tables/footnotes: none.
% Status: source-reviewed and compile-clean.
% Uncertainties: SOURCE ERRATA -- (11.1.19) uses the four-dimensional
%                determinant, and the prose prints A Delta s/T; corrected.

\subsection{Differential Equations for Stellar Structure}\label{sec:11.1}

We first set up the general-relativistic machinery for computing the pressure,
density, and gravitational fields within a spherically symmetric static star.
The metric will be taken in the ``standard'' form discussed in
Section~\ref{sec:8.1}:
\begin{equation}
  g_{rr}=A(r),\qquad
  g_{\theta\theta}=r^2,\qquad
  g_{\varphi\varphi}=r^2\sin^2\theta,\qquad
  g_{tt}=-B(r),\qquad
  g_{\mu\nu}=0\quad(\mu\ne\nu).
  \tag{11.1.1}\label{eq:11.1.1}
\end{equation}
The energy-momentum tensor is assumed to be that for a perfect fluid (see
Section~\ref{sec:5.4}):
\begin{equation}
  T_{\mu\nu}
  =p\,g_{\mu\nu}+(\rho+p)U_\mu U_\nu,
  \tag{11.1.2}\label{eq:11.1.2}
\end{equation}
with \(p\) the proper pressure, \(\rho\) the proper total energy density, and
\(U^\mu\) the velocity four-vector, defined so that
\begin{equation}
  g^{\mu\nu}U_\mu U_\nu=-1.
  \tag{11.1.3}\label{eq:11.1.3}
\end{equation}
Since the fluid is at rest, we take
\begin{equation}
  U_r=U_\theta=U_\varphi=0,\qquad
  U_t=-(-g^{tt})^{-1/2}=-\sqrt{B(r)}.
  \tag{11.1.4}\label{eq:11.1.4}
\end{equation}
Our assumptions of time independence and spherical symmetry imply that \(p\)
and \(\rho\) are functions only of the radial coordinate \(r\).

By making use of Eqs.~\eqref{eq:11.1.1}--\eqref{eq:11.1.4} and the Ricci
tensor components given by Eq.~\eqref{eq:8.1.13}, we find that the Einstein
equations read
% MODERNIZATION CHECK: Weinberg's Ricci tensor has the opposite sign.
% Both sides of (11.1.5)--(11.1.7) have been reversed for the binding
% curvature convention; the physical stellar equations below are unchanged.
\begin{align}
  R_{rr}
  &=
  -\frac{B''}{2B}
  +\frac{B'}{4B}\left(\frac{A'}A+\frac{B'}B\right)
  +\frac{A'}{rA}
  =4\pi G(\rho-p)A,
  \tag{11.1.5}\label{eq:11.1.5}\\
  R_{\theta\theta}
  &=
  1+\frac{r}{2A}\left(\frac{A'}A-\frac{B'}B\right)-\frac1A
  =4\pi G(\rho-p)r^2,
  \tag{11.1.6}\label{eq:11.1.6}\\
  R_{tt}
  &=
  \frac{B''}{2A}
  -\frac{B'}{4A}\left(\frac{A'}A+\frac{B'}B\right)
  +\frac{B'}{rA}
  =4\pi G(\rho+3p)B.
  \tag{11.1.7}\label{eq:11.1.7}
\end{align}
A prime denotes \(d/\dd r\). (We do not need to write down the equation for
\(R_{\varphi\varphi}\), which is identical to that for
\(R_{\theta\theta}\), or the equations for off-diagonal elements of
\(R_{\mu\nu}\), which simply say that zero equals zero.) In addition, we may
recall the equation for hydrostatic equilibrium,
\begin{equation}
  \frac{B'}B=-\frac{2p'}{p+\rho}.
  \tag{11.1.8}\label{eq:11.1.8}
\end{equation}

Our first step in solving these equations is to derive an equation for
\(A(r)\) alone, by forming the quantity
\begin{equation}
  \frac{R_{rr}}{2A}
  +\frac{R_{\theta\theta}}{r^2}
  +\frac{R_{tt}}{2B}
  =
  \frac{A'}{rA^2}+\frac1{r^2}-\frac1{Ar^2}
  =8\pi G\rho.
  \tag{11.1.9}\label{eq:11.1.9}
\end{equation}
This equation can be written
\begin{equation}
  \left(\frac rA\right)'=1-8\pi G\rho r^2.
  \tag{11.1.10}\label{eq:11.1.10}
\end{equation}
The solution with \(A(0)\) finite is
\begin{equation}
  A(r)=\left[1-\frac{2G\mathcal{M}(r)}r\right]^{-1},
  \tag{11.1.11}\label{eq:11.1.11}
\end{equation}
where
\begin{equation}
  \mathcal{M}(r)=\int_0^r4\pi r'^2\rho(r')\,\dd r'.
  \tag{11.1.12}\label{eq:11.1.12}
\end{equation}
We can now use Eqs.~\eqref{eq:11.1.11} and \eqref{eq:11.1.8} to eliminate
the gravitational fields \(A(r)\), \(B(r)\) from
Eq.~\eqref{eq:11.1.6}. We rewrite the result as
\begin{equation}
  -r^2p'(r)
  =
  G\mathcal{M}(r)\rho(r)
  \left[1+\frac{p(r)}{\rho(r)}\right]
  \left[1+\frac{4\pi r^3p(r)}{\mathcal{M}(r)}\right]
  \left[1-\frac{2G\mathcal{M}(r)}r\right]^{-1}.
  \tag{11.1.13}\label{eq:11.1.13}
\end{equation}
The reader may recognize this differential equation as the fundamental
equation of Newtonian astrophysics (see Section~\ref{sec:11.3}), with
general-relativistic corrections supplied by the last three factors.

We are primarily concerned in this chapter with stars that are isentropic,
that is, in which the entropy per nucleon \(s\) does not vary throughout the
star. This is the case for two very different kinds of star:

\emph{(A) Stars at Absolute Zero.} When a star exhausts its thermonuclear fuel
it can become a white dwarf (Section~\ref{sec:11.3}), or a neutron star
(Section~\ref{sec:11.4}), in which the temperature is essentially at absolute
zero. According to Nernst's theorem, the entropy per nucleon will then be zero
throughout the star.

\emph{(B) Stars in Convective Equilibrium.} If the most efficient mechanism
for energy transfer within the star is convection, then in equilibrium the
entropy per nucleon must be nearly constant throughout the star, because
otherwise a small element of fluid containing \(A\) nucleons could gain or
lose an energy \(A T\,\Delta s\) when transported from one part
of the star to another, and convection would therefore disturb the energy
distribution. The supermassive ``stars'' discussed in
Section~\ref{sec:11.5} are generally presumed to be in convective equilibrium.
% SOURCE ERRATUM: the scan prints A Delta s/T; thermodynamics gives A T Delta s.

We also assume that the stars we consider have a chemical composition that is
constant throughout.

The importance of the preceding assumptions lies in the fact that the pressure
\(p\) may in general be expressed as a function of the density \(\rho\), the
entropy per nucleon \(s\), and the chemical composition. Hence, with \(s\) and
the chemical composition constant throughout the star, \(p(r)\) may be
regarded as a function of \(\rho(r)\) alone, with no explicit dependence on
\(r\).

Given \(p(r)\) as a function \(p(\rho(r))\), we now formulate our problem as a
pair of first-order differential equations for \(\rho(r)\) and
\(\mathcal{M}(r)\). One of these is Eq.~\eqref{eq:11.1.13}; the other is the
derivative of Eq.~\eqref{eq:11.1.12}:
\begin{equation}
  \mathcal{M}'(r)=4\pi r^2\rho(r).
  \tag{11.1.14}\label{eq:11.1.14}
\end{equation}
In addition, Eq.~\eqref{eq:11.1.12} provides an initial condition:
\begin{equation}
  \mathcal{M}(0)=0.
  \tag{11.1.15}\label{eq:11.1.15}
\end{equation}
Equations~\eqref{eq:11.1.13}--\eqref{eq:11.1.15}, together with an equation
of state giving \(p(\rho)\), serve to determine \(\rho(r)\),
\(\mathcal{M}(r)\), \(p(r)\), and so on, throughout the star, once we specify
the other initial condition, that is, the value of \(\rho(0)\). The
differential equations \eqref{eq:11.1.13} and \eqref{eq:11.1.14} must be
integrated out from the center of the star, until \(p(\rho(r))\) drops to zero
at some point \(r=R\), which we then interpret as the radius of the particular
star with central density \(\rho(0)\).

Let us return to the problem of calculating the metric. Once we compute
\(\rho(r)\), \(\mathcal{M}(r)\), and \(p(r)\), we can immediately obtain
\(A(r)\) from Eq.~\eqref{eq:11.1.11}; to find \(B(r)\) we use
Eq.~\eqref{eq:11.1.13} to rewrite Eq.~\eqref{eq:11.1.8} as
\[
  \frac{B'}B
  =
  \frac{2G}{r^2}\left[\mathcal{M}+4\pi r^3p\right]
  \left[1-\frac{2G\mathcal{M}}r\right]^{-1}.
\]
The solution with \(B(\infty)=1\) is
\begin{equation}
  B(r)
  =
  \exp\left\{
    -\int_r^\infty
    \frac{2G}{r'^2}
    \left[\mathcal{M}(r')+4\pi r'^3p(r')\right]
    \left[1-\frac{2G\mathcal{M}(r')}{r'}\right]^{-1}
    \dd r'
  \right\}.
  \tag{11.1.16}\label{eq:11.1.16}
\end{equation}
Our solution is now complete. (Incidentally, we did not need to use
Eqs.~\eqref{eq:11.1.5} and \eqref{eq:11.1.7} for \(R_{rr}\) and \(R_{tt}\)
separately, because these equations follow from Eqs.~\eqref{eq:11.1.6},
\eqref{eq:11.1.8}, and \eqref{eq:11.1.9}, which were used in our
calculation. This should not be surprising, because Eq.~\eqref{eq:11.1.8},
which is really just the equation for momentum conservation, follows from the
Einstein equations \eqref{eq:11.1.5}--\eqref{eq:11.1.7} via the Bianchi
identities.)

Outside the star, \(\rho(r)\) and \(p(r)\) vanish, and \(\mathcal{M}(r)\) is
the constant \(\mathcal{M}(R)\). Thus Eqs.~\eqref{eq:11.1.11} and
\eqref{eq:11.1.16} give
\begin{equation}
  B(r)=A^{-1}(r)=1-\frac{2G\mathcal{M}(R)}r
  \qquad\text{for }r\ge R.
  \tag{11.1.17}\label{eq:11.1.17}
\end{equation}
The discussion of Section~\ref{sec:8.2} shows that the constant
\(\mathcal{M}(R)\) that appears in the asymptotic gravitational field
\eqref{eq:11.1.17} must equal the mass \(M\) of the star, defined as the total
energy of the star and its gravitational field, that is,
\begin{equation}
  M=\mathcal{M}(R)=\int_0^R4\pi r^2\rho(r)\,\dd r.
  \tag{11.1.18}\label{eq:11.1.18}
\end{equation}
Thus Eq.~\eqref{eq:11.1.17} is just the familiar exterior Schwarzschild
solution.

It may appear paradoxical that \(M\), which must include the energy of the
gravitational field, is given in Eq.~\eqref{eq:11.1.18} as the integral of
the energy density \(\rho(r)\) of matter (including radiation) alone. The
resolution is that Eq.~\eqref{eq:11.1.18} does \emph{not} say that \(M\) is
the total energy of the matter. The total material energy is not really well
defined, but it might be computed by splitting up the star into small volume
elements and adding up the energies of each element as measured in a locally
inertial reference frame; this would give the material energy as
% SOURCE ERRATUM: the scan uses the four-dimensional determinant; material
% energy on a static slice requires the induced proper three-volume measure.
\begin{equation}
  M_{\mathrm{matter}}
  =
  \int\sqrt{A(r)}\,\rho(r)\,r^2\sin\theta\,\dd r\,\dd\theta\,\dd\varphi
  =
  \int_0^R4\pi r^2\sqrt{A(r)}\,\rho(r)\,\dd r.
  \tag{11.1.19}\label{eq:11.1.19}
\end{equation}
The difference between Eqs.~\eqref{eq:11.1.18} and \eqref{eq:11.1.19} can
be regarded as the energy of the gravitational field. However, this
decomposition is not particularly useful, and will not be employed here.

It is more informative to compare Eq.~\eqref{eq:11.1.18} with the energy
\(M_0\) that the matter of the star would have if dispersed to infinity. This
is simply
\begin{equation}
  M_0=m_NN,
  \tag{11.1.20}\label{eq:11.1.20}
\end{equation}
where \(m_N=1.66\times10^{-24}\ \mathrm{g}\) is the rest mass of a nucleon and
\(N\) is the number of nucleons in the star. The nucleon number is given by
\begin{equation}
  N
  =
  \int\sqrt{-g}\,J_N^0\,\dd r\,\dd\theta\,\dd\varphi
  =
  \int_0^R4\pi r^2\sqrt{A(r)B(r)}\,J_N^0(r)\,\dd r,
  \tag{11.1.21}\label{eq:11.1.21}
\end{equation}
where \(J_N^\mu\) is the conserved nucleon-number current. It is convenient
to express \(J_N^0\) in terms of the proper nucleon-number density \(n\), that
is, the nucleon-number density measured in a locally inertial reference frame
at rest in the star, which is
\begin{equation}
  n=-U_\mu J_N^\mu=\sqrt{B}\,J_N^0.
  \tag{11.1.22}\label{eq:11.1.22}
\end{equation}
(See Eq.~\eqref{eq:11.1.4}, and recall that in a locally inertial coordinate
frame \(U_0=-1\).) Equation~\eqref{eq:11.1.21} then becomes
\begin{equation}
  \begin{split}
  N
  &=
  \int_0^R4\pi r^2\sqrt{A(r)}\,n(r)\,\dd r\\
  &=
  \int_0^R4\pi r^2
  \left[1-\frac{2G\mathcal{M}(r)}r\right]^{-1/2}
  n(r)\,\dd r.
  \end{split}
  \tag{11.1.23}\label{eq:11.1.23}
\end{equation}
The proper number density \(n(r)\) is in general a function of the proper
density \(\rho(r)\), the chemical composition, and the entropy per nucleon
\(s\), so \(n(r)\) and \(N\) are fixed for a star with a given constant \(s\)
and chemical composition, once we choose \(\rho(0)\). The internal energy of
the star is now given by
\begin{equation}
  E\equiv M-m_NN.
  \tag{11.1.24}\label{eq:11.1.24}
\end{equation}
We can also define a proper internal material energy density as
\begin{equation}
  e(r)\equiv\rho(r)-m_Nn(r)
  \tag{11.1.25}\label{eq:11.1.25}
\end{equation}
and write Eq.~\eqref{eq:11.1.24} as
\begin{equation}
  E=T+V,
  \tag{11.1.26}\label{eq:11.1.26}
\end{equation}
where \(T\) and \(V\) are the thermal and gravitational energies,
respectively, of the star:
\begin{align}
  T
  &\equiv
  \int_0^R4\pi r^2
  \left[1-\frac{2G\mathcal{M}(r)}r\right]^{-1/2}
  e(r)\,\dd r,
  \tag{11.1.27}\label{eq:11.1.27}\\
  V
  &\equiv
  \int_0^R4\pi r^2
  \left\{
    1-\left[1-\frac{2G\mathcal{M}(r)}r\right]^{-1/2}
  \right\}
  \rho(r)\,\dd r.
  \tag{11.1.28}\label{eq:11.1.28}
\end{align}
Expanding the square roots gives
\begin{align}
  T
  &=
  \int_0^R4\pi r^2
  \left\{1+\frac{G\mathcal{M}(r)}r+\cdots\right\}
  e(r)\,\dd r,
  \tag{11.1.29}\label{eq:11.1.29}\\
  V
  &=
  -\int_0^R4\pi r^2
  \left\{
    \frac{G\mathcal{M}(r)}r
    +\frac{3G^2\mathcal{M}^2(r)}{2r^2}
    +\cdots
  \right\}\rho(r)\,\dd r.
  \tag{11.1.30}\label{eq:11.1.30}
\end{align}
The first terms in \(T\) and \(V\) are recognizable as the Newtonian values
for the thermal and gravitational energies of the star; in particular, note
that the first term in \(V\) may be written
\begin{equation}
  \begin{aligned}
  -G\int_0^R4\pi r\mathcal{M}(r)\rho(r)\,\dd r
  &=-\frac G2\int_0^R\frac1r\,\dd\!\left(\mathcal{M}^2(r)\right)\\
  &=-\frac{GM^2}{2R}
    -\frac G2\int_0^R\frac{\mathcal{M}^2(r)}{r^2}\,\dd r\\
  &=\frac{\phi(R)\mathcal{M}(R)}2
    -\frac12\int_0^R\mathcal{M}(r)\,\dd\phi(r)\\
  &=\frac12\int_0^R\phi(r)\,\dd\mathcal{M}(r),
  \end{aligned}
  \tag{11.1.31}\label{eq:11.1.31}
\end{equation}
where \(\phi\) is the Newtonian potential, given inside the star by
\[
  \phi(r)=-\frac{GM}{R}
  -G\int_r^R\frac{\mathcal{M}(r')}{r'^2}\,\dd r'.
\]
The higher terms in \(T\) and \(V\) are discussed in
Section~\ref{sec:11.5}.

To repeat our main conclusion: once we specify that a star has a definite
uniform entropy per nucleon and chemical composition, all properties of the
star, including \(\rho(r)\), \(p(r)\), \(n(r)\), \(e(r)\), \(M\), \(N\), and
\(E\), are determined as functions of the central density \(\rho(0)\). This is
not the case for ordinary stars like the sun, in which the entropy distribution
is not uniform and has to be determined from the equations of radiative
equilibrium. However, the considerations of this section do provide an
adequate basis for the study of the exotic structures discussed in this
chapter.
